\par\noindent\begin{minipage}{\linewidth}
\begin{summarybox}
\textbf{核心问题：}已知顶点坐标，怎么最快求内心？已知向量等式，怎么识别内心？

\medskip
\textbf{一、最少记忆量：只记住对应关系}
\[
\underbrace{BC}_{a}\longleftrightarrow A,\qquad
\underbrace{CA}_{b}\longleftrightarrow B,\qquad
\underbrace{AB}_{c}\longleftrightarrow C.
\]
\[
\boxed{a\vec{IA}+b\vec{IB}+c\vec{IC}=\vec0},\quad
\boxed{I=\frac{aA+bB+cC}{a+b+c}\ \text{（表示坐标逐项计算）}}.
\]
第二个框里的 $A,B,C,I$ 仅是\textbf{坐标的简写}；正式向量写法是
\[
\vec{OI}=\frac{a\vec{OA}+b\vec{OB}+c\vec{OC}}{a+b+c}.
\]

\textbf{二、看见什么，用什么}
\begin{enumerate}
\item \textbf{已知三顶点求内心：}先算三条对边，再做两次按边长乘坐标的加法。
\item \textbf{向量等式系数是对边长：}检查次序，直接识别内心；若系数是同一倍数，先约去倍数。
\item \textbf{只求角平分线交点：}先用 $BD:DC=c:b$，不必先求整个内心。
\item \textbf{只求角平分线方向：}两条边的单位向量相加；求位置还要增加条件。
\item \textbf{问三个小三角形面积比：}同高，直接得到 $a:b:c$。
\end{enumerate}

\textbf{三、三种主要证明怎么选择}
\begin{itemize}
\item \textbf{想真正弄懂：}先读“初中面积和比例”证明，理解两次沿线段定位。
\item \textbf{考场追求书写少：}掌握面积定位的小结论后，两步从面积比得到向量式。
\item \textbf{常规且稳妥：}两次角平分线定理，再用向量分点公式。
\end{itemize}
文中另有\textbf{坐标距离独立验证}，用来确认结果没有靠图猜测。

\textbf{四、关联知识}

角平分线定理 $\to$ 线段比例 $\to$ 向量分点；
内切圆半径相等 $\to$ 三块面积比例 $\to$ 内心坐标；
重心公式的三个相同系数 $\leftrightarrow$ 内心公式的三个对边系数。

\medskip
\textbf{最后提醒：}不要靠“内心权重”四个字硬记。\textbf{每个系数来自该顶点对面的边长}，这是最不容易记错的根由。
\end{summarybox}
\end{minipage}\par
